% Source: https://artofproblemsolving.com/community/c3h2373945_what_is_a_function
% Stable identity: aops-topic-2373945-q01
\begin{enumerate}
\item[(a)] \textbf{Function.} Every real $x$ has exactly one square $x^2$. (Different inputs, such as $2$ and $-2$, may share an output; that is allowed.)
\item[(b)] \textbf{Not a function.} The rule does not give exactly one output: for $x=4$, both $y=2$ and $y=-2$ satisfy $y^2=4$. (Choosing the nonnegative one, $x\mapsto\sqrt x$, does give a function.)
\item[(c)] \textbf{Function.} Each of $1,2,3$ gets exactly one output; two inputs sharing the output $4$ is allowed. (It would also be allowed for some elements of $Y$ to be unused.)
\item[(d)] \textbf{Not a function.} The input $1$ is assigned two outputs, $4$ and $5$.
\item[(e)] \textbf{Not a function} from all real numbers: $x=0$ gets no output. (It is a function from the nonzero real numbers to the real numbers.)
\end{enumerate}
So exactly $\boxed{\text{(a) and (c)}}$ define functions. For a rule given by a graph in the plane, this is the vertical line test: each vertical line $x=c$ with $c$ in $X$ must meet the graph exactly once.
